\section{A secondary improvement: twelve positive unknowns at degree 938}
\label{sec:positive-majorant}

This section uses the complete source-31 coefficient tuple already bundled
with this report, with port aliases
\(\beta=\texttt{beta},K_0=\texttt{K},M_C=\texttt{MC},
M_F=\texttt{MF\_source}\). The last mask is shifted. It does not use the
Sun coefficient or the new exponent code. The inherited normalized
blank-K215 compiler interface is the one proved in source 31's
\texttt{coefficients/COEFFICIENT\_PROOF.md}; source 34's positive reductions
are explicitly uniform in such complete tuples. The new step is a
separator of formal degree nine, replacing the previous degree-23
separator. All symbols in this section are local.

The twelve supplied unknowns are
\[
J,F,\alpha,z,h,i,T,s,w,\eta,\zeta,t\in\Z_{>0}.
\]
All other symbols below are polynomial abbreviations. Set
\begin{align*}
q&=\beta J+1,&\omega&=w+q,&C&=q-F-\alpha,\\
X&=\omega q,&Y&=sq^3,&E&=XY,\\
k&=\eta+\zeta,&c&=kY+\eta,&a&=E+Y,\\
H&=4a+3,&\Delta&=a^2+H,&\ell&=EkY,\\
A&=1+\ell(\ell+k),&G&=1+\Delta c^2,&b_0&=X+ac,\\
I&=ic^2,&D&=1+\Delta I^2,&Q&=(\Delta I)^2,\\
d&=Q-1,&v&=cT,&b&=c+RD.
\end{align*}
Here the abbreviation \(R\), and the two transport expressions, are
\begin{align*}
R&=((q-1)(q-F)+\alpha)(q^2-1)+(M_C+qM_F)J,\\
N_k&=k-hE-R,\qquad
N_t=(K_0+\omega)C+q-F-z(q-1).
\end{align*}
The new majorant and translated carrier are
\begin{gather*}
\boxed{M=k(ik^3+T)^2},\qquad \delta=(a+2)c-\ell,\\
\sigma=b+M\delta+v(a+2)I,\qquad
L=HM+1,\qquad W=b_0+H(t+\sigma).
\end{gather*}
Write \(e=Q(Dv^2+b^2)-1\) and \(B_1=2Qvb\). The polynomial norm is
\begin{equation}
N=\prod_{\substack{\tau=\pm\sqrt A,\;g=\pm\sqrt G\\f=\pm\sqrt D}}
\bigl[d(W-HM\tau-Lg)^2-H^2e+H^2B_1f\bigr].
\label{eq:positive-newnorm}
\end{equation}
This is an integer polynomial; the formal sign product introduces no
quantified unknowns. It is valid also for square, zero, or dependent
radicands. The scalar circuit below and the delivered JSON files give
a literal evaluation using only integer additions and multiplications.
The two presentations are
\begin{equation}
\boxed{N_kN_t(1+N)-1=0},\qquad
\boxed{N_k=1,\quad N_t=1,\quad N=0}.
\label{eq:positive-presentations}
\end{equation}

For completeness, one scalar evaluation of the norm is
\begin{align*}
n&=(HM)^2A,&d_n&=dn,&H_Q&=H^2Q,\\
H_P&=H_QDv^2,&H_U&=H_Qb^2,&H_e&=H_P+H_U-H^2,\\
j&=4H_PH_U,&\mathcal T&=2d(W^2+L^2G),&
\mathcal K&=[d(W^2-L^2G)]^2,\\
a_t&=\mathcal T-2(d_n+H_e),&
b_t&=(d_n-H_e)^2-j-\mathcal K,\\
\bar v&=16d_nj-a_t(a_t\mathcal T+2b_t),&
u&=b_t^2-a_t^2\mathcal K,\\
N&=\mathcal K\bar v^2-u(\mathcal T\bar v-u).
\end{align*}
These are the sourced trace--remainder identities with the new inputs.

\begin{theorem}
For the complete fixed coefficient tuple specified above, ZFC proves
that absence of positive integer solutions to either presentation in
\eqref{eq:positive-presentations} is equivalent to \(\operatorname{Con}(\mathrm{ZFC})\).
The single polynomial has exactly twelve supplied unknowns, degree 938,
and a 122-operation evaluator. The system has three equations, maximum
residual degree 929, and 117 operations. Binary additions, subtractions,
and multiplications by numerals are all charged; reuse and system
comparisons are free.
\end{theorem}

\paragraph{Inherited interface, stated explicitly.}
For the authentic tuple, \(N_k=N_t=\epsilon\),
\(\epsilon\in\{1,-1\}\), implies before recovering any root that
\begin{gather*}
q\ge16,\quad Y\ge q^3,\quad E=XY>R+2,\quad
3q<R<q^4-2,\\
k=R+\epsilon+hE>E,\quad kY<c<k(Y+1),\quad
\sqrt G-\sqrt A>kY^2.
\end{gather*}
The recovered four normalized Pell equations, with integral roots and
\(g\equiv b_0\pmod H\), imply a blank K215 halt; the signed-quotient
extension restores \(\epsilon=1\). Conversely a halt admits canonical
positive roots \(g,\tau,f,y,\gamma\), with \(X=2^R\).
These are the inherited compiler assertions, not claims established by
finite tests in the present report. The coordinate shift is complete:
\(2^R/q>q\), so the original \(w_{\rm old}\) can be replaced by the
positive integer \(w=w_{\rm old}-q\). This is the shift proved in source 37.

\paragraph{The new majorant estimate.}
The shifted coordinate gives \(X\ge q(q+1)>q^2\), hence \(R<X^2\).
The bounds
\(c<(17/16)kY\) and \(a+2<(17/16)E\) follow at once.
Crucially,
\begin{equation}
\frac{R(a+2)^2c^4}{k^8Y^2}<1.
\label{eq:new-majorant-key}
\end{equation}
Indeed the left side equals
\[
\frac R{X^2}\left(\frac{a+2}{E}\right)^2
\left(\frac c{kY}\right)^4\left(\frac Ek\right)^4
\le\left[\frac q{q+1}(1+q^{-2}+2q^{-5})(1+q^{-3})^2\right]^2<1.
\]
For the last inequality, set \(x=q^{-2}\). The two parentheses have
product at most \((1+2x)(1+x)^2<1+5x<1+1/q\), since \(q\ge16\).

Let \(Z=v\sqrt D+b\). Every complex auxiliary root satisfying
\(dy^2=Q(vf-b)^2-1\), \(f^2=D\), obeys the sharp estimate
\begin{equation}
|y|<Z+v+1.
\label{eq:positive-sharp-y}
\end{equation}
To see this, \(\Delta>R+2\), \(I^2\ge c^4\) give
\(Q-b=\Delta I^2(\Delta-R)-c-R>0\), so \(b\le d\), and
\(\sqrt D<\Delta I<d\). Therefore
\[
|y|^2\le\frac{QZ^2+1}d
<\left[Z\left(1+\frac1d\right)\right]^2,
\qquad Z/d<v+1.
\]
The strict middle inequality has numerator \(QZ^2-d>0\).

Since \(D<(a+2)^2I^2\), equation \eqref{eq:new-majorant-key} yields
\(b/(kY^2)<i^2k^7+1\). Also
\[
\frac{v\sqrt D}{kY^2}
<(17/16)^4EYiTk^2
<\frac{(17/16)^4}{X}iTk^4<iTk^4.
\]
Using \(v/(kY^2)<T\) and \(1/(kY^2)<1\), we obtain
\[
\frac{Z+v+1}{kY^2}<i^2k^7+iTk^4+T+2
<i^2k^7+2iTk^4+kT^2=M.
\]
Thus \(|y|<MkY^2<M(\sqrt G-\sqrt A)\).
Moreover \(\sqrt A>\ell=XkY^2\), so
\(M\sqrt A>XZ>2Z\). The latter will ensure positivity in the converse.

\paragraph{Soundness.}
The scalar norm satisfies \(N\equiv\square\pmod4\): write
\(\mathcal T=2x_0\) and
\(\mathcal K=x_0^2-4d^2W^2L^2G\) in the final formula. Thus a zero of
the product in \eqref{eq:positive-presentations} forces \(N=0\) and
\(N_k=N_t=\epsilon\in\{1,-1\}\). A vanishing factor gives coherent
algebraic roots with
\[
x=(W-Lg)/H=M\tau+y,\quad \tau^2=A,\quad g^2=G,\quad f^2=D,
\quad dy^2=Q(vf-b)^2-1.
\]
Every conjugate of \(x\) has magnitude less than \(M\sqrt G\).
If \(g\) were irrational, an automorphism reversing it would produce
two such \(x\)'s at distance \(2L\sqrt G/H>2M\sqrt G\), impossible.
Hence \(g\) is rational and integral, and \(x\) is rational. A sign
change of an irrational \(\tau\) would require two auxiliary roots of
magnitude less than \(M\sqrt A\) to compensate a difference
\(2M\sqrt A\). Thus \(\tau\) is integral, \(y\) is rational, and
\(f=(e-dy^2)/B_1\) is rational and integral.

The sourced rational Pell saturation lemma now applies to
\[
(\Delta I)^2(vf-b)^2-((\Delta I)^2-1)y^2=1
\]
and makes \(y\), hence \(x\), integral. This lemma says that for
integers \(S\ge2,V\), a rational \(y\) satisfying
\(S^2V^2-(S^2-1)y^2=1\) is integral. Source 24 proves it by reducing
the denominator to a common Pell unit and using divisibility of its
cosine sequence.

Since \(\delta,\sigma>0\), \(W>0\); a negative \(g\) would contradict
the bound on \(x\). Reducing \(W-Lg=Hx\) modulo \(H\) gives
\(g\equiv b_0\pmod H\). Its quotient \(\gamma\) is positive because
\(g>(a+1)c>b_0\). The normalized compiler interface now supplies a blank
K215 halt; its ordinary strong coordinate is \(i_{\rm ord}=\Delta i\).
No independent sign changes of dependent radicals were assumed.

\paragraph{Completeness.}
From a canonical positive compiler witness and the complete \(w\)-shift,
put
\[
t=\gamma+M(g-\delta)+M\tau+y-b-v(a+2)I.
\]
We have \(0<\delta<ac<g\), because \(\ell/c>E/2>2\).
Also \((a+2)I<2\sqrt D\) and
\(M\tau>2(v\sqrt D+b)>b+v(a+2)I\). Therefore \(t\) is a positive
integer. Its definition makes \(W=Lg+H(M\tau+y)\), so one norm factor
vanishes and \(N_k=N_t=1\). Composing both directions with the calibrated
K215 refutation theorem, and negating existential solvability, proves the
stated ZFC equivalence.

\paragraph{Exact degree.}
Scale the twelve variables by a formal parameter. The degrees of
\(E,H,c,I,D,Q,b,M,v\) are respectively
\(6,6,5,11,34,46,38,9,6\).
The positive-leading roots \(g,\tau,f\) have degrees \(11,11,17\), with
\[
g-\tau=\delta-k/2+O(\lambda^{-1}),\qquad
(a+2)I-f=I/(2(a+2))+\text{lower-degree terms}.
\]
The four positive-\(f\) factors in \eqref{eq:positive-newnorm} have degree
119. In the negative-\(f\) factors the leading auxiliary terms cancel.
Their remaining leading terms have degrees 116,116,114,107: the first
three come from the main radical signs, and the last from
\(Hv((a+2)I-f)\), of degree 17. All competing squared terms have degree
at most 104. Thus \(\deg N=4(119)+2(116)+114+107=929\).

More explicitly, with \(q_0=\beta J\), \(k_0=\eta+\zeta\),
\(\omega_0=w+q_0\), the nonzero homogeneous leader is
\[
N_{929}=-2^{46}i^{43}k_0^{103}T^5s^{152}
\omega_0^{69}q_0^{549}(q_0-F)^8.
\]
This is nonzero for every authentic \(\beta>0\). Since
\(\deg N_k=7\) and \(\deg N_t=2\), the product has exact degree 938.
The new product improves source 34's degree-1022, 124-operation form in
both metrics; its system improves the corresponding degree-1013,
120-operation system. The existing degree-1081 evaluator at 108 operations
retains its operation-count advantage. There is no variable-count or
global-optimality claim.

The delivered checker compares 123 independent formal-radical products
with the scalar circuit and verifies six complete univariate degree
specializations, including three with the actual coefficient tuple.
The generic runs also check the displayed leading coefficient. The full
proof details are in \texttt{positive/POSITIVE\_PROOF.md}. These finite
checks support the algebra and do not replace the inherited compiler
theorems or claim a new end-to-end formalization.
